% Compile: pdflatex -enable-write18 top.tex
% Include graphics from Hefferon book: rename file to have .eps extension
% VS code: re-enable auto build: latex-workshop.latex.autoBuild.run
%
\documentclass[9pt,aspectratio=169,xcolor=dvipsnames]{beamer}

\usetheme{Singapore}
\usecolortheme{seagull}

\definecolor{TitleBlue}{RGB}{0, 90, 156}
\setbeamercolor{frametitle}{fg=TitleBlue}
\setbeamercolor{structure}{fg=TitleBlue}

\setbeamerfont{frametitle}{series=\bfseries,parent=structure}
\setbeamerfont{title}{size=\LARGE}
% \usetheme[background=light,sectionstyle=style2,slidestyle=cyber,numbering=fraction,headingcolor=theme]{trigon}

\usepackage{mylinsys}
\usepackage{graphicx}
\usepackage{grffile}
\usepackage{amsfonts}
\usepackage{ucs}
\usepackage[utf8x]{inputenc}
\usepackage{physics}  % \va \vb
\usepackage{bm}
\usepackage{upgreek}
\usepackage{svg}
\usepackage{algpseudocode}
\usepackage{marvosym}
\usepackage{hyperref}
\usepackage{listings}

\renewenvironment{boxed}
{
    \par\addvspace{0.5em}\noindent

    \definecolor{blizzard}{rgb}{0.9375, 0.96875, 0.99609375}
    \setbeamercolor{postit}{fg=black,bg=blizzard}

    \begin{beamercolorbox}[sep=0.25em,rounded=true,wd=\textwidth]{postit}
        }
        {
    \end{beamercolorbox}

    \par\addvspace{0.5em}
}

\newenvironment{conceptcheckbox}
{
    \par\addvspace{0.75em}\noindent

    \definecolor{lightblue}{RGB}{125, 249, 255}
    \setbeamercolor{conceptcheckboxc}{fg=black, bg=lightblue}

    \begin{beamercolorbox}[sep=0.4em,rounded=true,wd=\textwidth]{conceptcheckboxc}

        \begin{columns}[onlytextwidth]
            \begin{column}{0.15\textwidth}
                \hspace*{2mm}
                \includegraphics[width=0.7\textwidth]{fig/conceptcheck.eps}
            \end{column}
            \begin{column}{0.85\textwidth}
                }
                {
            \end{column}
        \end{columns}

    \end{beamercolorbox}

    \par\addvspace{0.75em}
}


\renewcommand{\v}[1]{\bm{#1}} % vector
\newcommand{\m}[1]{\bm{#1}} % matrix
\renewcommand{\d}[1]{\mathop{d#1}} % differential
\renewcommand{\ev}[1]{\mathcal{#1}} % event (probability)
\newcommand{\E}[1]{\operatorname{E}\left[#1\right]} % expected value
\newcommand{\const}{\mathrm{const}} % matrix
\renewcommand{\mid}{\,|\,} % conditional "bar" (probability)
\newcommand{\conceptcheck}[2]{
    \begin{frame}{}
        \begin{conceptcheckbox}
            #1
        \end{conceptcheckbox}
    %\vspace*{2mm}
    %\textbf{Answer:}\\[1mm]
    %#2
    \end{frame}
}

\begin{document}

\begin{frame}
    % manual titlepage
    \begin{columns}[onlytextwidth]
        \begin{column}{0.4\textwidth}
            \vspace*{10mm}
            \includegraphics[width=1.4\textwidth]{fig/titleimage.eps}
            {\tiny Image by \href{https://pixabay.com/users/julia_s-4109630/}{Julia S} from
                \href{https://pixabay.com}{Pixabay}}
        \end{column}
        \begin{column}{0.6\textwidth}
            \flushright
            {A Refresher On}\\[1mm]
            {\Huge \textbf{Machine Learning Maths}}\\[15mm]
            {\Large CAS Machine Learning}\\[1mm]
            {\Large Module 1: Introduction}\\[15mm]
            Michael Calonder\\[1mm]
            Hochschule Luzern HSLU\\[1mm]
            {March 15 \& 16, 2024}\\
        \end{column}
    \end{columns}
\end{frame}

\begin{frame}{Outline}
    \tableofcontents
\end{frame}

\section{Part 1: Linear Algebra}

\subsection{Vectors and Matrices}

\begin{frame}{Linear Algebra}
    \framesubtitle{Introduction}
    What is linear algebra and where is it used?
    \begin{itemize}
        \item Branch of mathematics concerned with linear equations such as
              $$a_1 x_1 + a_2 x_2 + \dots + a_n x_n = b$$
        \item Of central important to most areas of science, mathematics, and engineering: ranging from differential equations, to modeling electrical circuits and control systems that fly aircraft, to econometry and genetics
        \item Not just ``linear'': Non-linear systems may well be treated through a series of linear operations, e.g. first-order Taylor expansion
        \item At the heart of most machine learning algorithms
        \item Modern computers can process them very efficiently, still active research (DeepMind's AlphaTensor)
    \end{itemize}
\end{frame}

\begin{frame}{Vectors}
    Properties:
    \begin{itemize}
        \item A vector is an object in $\mathbb{R}^n$ that is comprised of a magnitude and a direction
        \item Often visualized as an ``arrow'' for $n=2$ or $n=3$\\
              \begin{center}
                  \includegraphics{gr/mp/ch1.12}
              \end{center}
              Note that all above arrow represent the same vector
        \item Notation: $$\vb{x} := \mqty(x_1 \\ x_2 \\ \vdots \\ x_n)$$
    \end{itemize}
\end{frame}


\begin{frame}{Vectors}
    \framesubtitle{Linear Dependence}
    \begin{boxed}
        Given $\v x_1, \ldots \v x_m$ be vectors in $\mathbb{R}^n$.

        The set $\{ \v x_1, \ldots \v x_m \}$ of these vectors \emph{linearly independent} if
        $$c_1 \v x_1 + c_2 \v x_2 + \ldots + c_m \v x_m \quad \text{implies} \quad c_1 = c_2 = \ldots = c_m = 0$$
    \end{boxed}
    Likewise, any two vectors $\v x_i$ and $\v x_j$ are \emph{linearly dependent} if $\v x_1 = c \v x_2$
    for some real $c \neq 0$\\[3mm]
\end{frame}

\conceptcheck{
    Determine whether or not the two vectors $\mqty(6, 2, 8, 4)^\top$
    and $\mqty(21,  7, 28, 14)^\top$ are linearly dependent
}{
    They are linearly dependent, $3.5 \; \times$ the first vector yields the second
}

\begin{frame}{Vectors}
    \begin{columns}[onlytextwidth]
        \begin{column}{0.5\textwidth}
            Algebraic properties:
            \begin{itemize}
                \item Additive identity: $\forall \v{x} : \exists \v{0}$ s.t. $\v{x} + \v{0} = \v{x}$
                \item Multiplicative identity: $1\, \v{x} = \v{x}$
                \item Commutativity: $\vb{x} + \vb{y} = \vb{y} + \vb{x}$
                \item Associativity: $(\vb{x} + \vb{y}) + \vb{z} = \vb{y} + (\vb{x} + \vb{z})$
                \item Additive inverse: $\forall \v{v} : \exists \bar{\v{x}} := -\v{x}$  s.t. $\v{x} + \bar{\v{x}} = \v{x} + (-\v{x}) = \v{0}$
                \item Distributivity:
                      \begin{itemize}
                          \item $\forall \v{v}, r$: $r(\v{x} + \v{y}) = r \v{x} + r \v{y}$
                          \item $\forall \v{v}, r, s$: $(r+s) \v{x} = r \v{x} + s \v{x}$
                      \end{itemize}
            \end{itemize}
            For $\v{x}, \v{y}, \v{z} \in \mathbb{R}^n$ and $r$, $s$ scalars

        \end{column}
        \begin{column}{0.5\textwidth}
            \begin{center}
                \includegraphics[width=0.5\textwidth]{gr/mp/ch1.13} \\[7mm]
                \includegraphics[width=0.4\textwidth]{gr/mp/ch1.14} \\[7mm]
                \includegraphics[width=0.9\textwidth]{gr/mp/ch1.17}
            \end{center}
        \end{column}
    \end{columns}

\end{frame}

\begin{frame}{Matrices}
    Basic properties:
    \begin{itemize}
        \item A matrix is an object in $\mathbb{R}^{m \times n}$, i.e. a rectangular array of numbers.
        \item Central object in linear algebra, enable actual operations, e.g. addition or multiplication
        \item Widely used, even outside linear algebra e.g. in graph theory (adjacency matrices)
        \item Notation:
              $$ \m{A} := \begin{pmatrix}
                      a_{11} & a_{12} & \cdots & a_{1n} \\
                      a_{21} & a_{22} & \cdots & a_{2n} \\
                      \vdots & \vdots & \ddots & \vdots \\
                      a_{m1} & a_{m2} & \cdots & a_{mn}
                  \end{pmatrix}
              $$
        \item Note that for now we choose $a_{ij}$ to be in $\mathbb{R}$, however, $a_{ij} \in \mathbb{C}$ in general
        \item Note that algebraic properties and identities differ form those above ones for vectors (see below)
        \item Important special case: \emph{square} matrix, $m = n$
    \end{itemize}
\end{frame}


\begin{frame}{Matrix Multiplication (MM)}
    Matrix product $\m{C} = \m{A}\m{B} \in \mathbb{R}^{m \times p}$ with $\m{A} \in \mathbb{R}^{m \times n}$,
    $\m{B} \in \mathbb{R}^{n \times p}$,
    $$\m{A}=\begin{pmatrix}
            a_{11} & a_{12} & \cdots & a_{1n} \\
            a_{21} & a_{22} & \cdots & a_{2n} \\
            \vdots & \vdots & \ddots & \vdots \\
            a_{m1} & a_{m2} & \cdots & a_{mn} \\
        \end{pmatrix},\quad\m{B}=\begin{pmatrix}
            b_{11} & b_{12} & \cdots & b_{1p} \\
            b_{21} & b_{22} & \cdots & b_{2p} \\
            \vdots & \vdots & \ddots & \vdots \\
            b_{n1} & b_{n2} & \cdots & b_{np} \\
        \end{pmatrix},\quad\m{C}=\begin{pmatrix}
            c_{11} & c_{12} & \cdots & c_{1p} \\
            c_{21} & c_{22} & \cdots & c_{2p} \\
            \vdots & \vdots & \ddots & \vdots \\
            c_{m1} & c_{m2} & \cdots & c_{mp} \\
        \end{pmatrix}
    $$
    is defined s.t.
    $$c_{ij} = a_{i1}b_{1j} + a_{i2}b_{2j} +\cdots + a_{in}b_{nj}= \sum_{k=1}^n a_{ik}b_{kj}$$

    For instance,
    $$
        \begin{pmatrix}
            a_{11} & a_{12} \\
            \cdot  & \cdot  \\
            a_{31} & a_{32} \\
            \cdot  & \cdot  \\
        \end{pmatrix}
        \begin{pmatrix}
            \cdot & b_{12} & b_{13} \\
            \cdot & b_{22} & b_{23} \\
        \end{pmatrix}
        = \begin{pmatrix}
            \cdot & c_{12} & \cdot  \\
            \cdot & \cdot  & \cdot  \\
            \cdot & \cdot  & c_{33} \\
            \cdot & \cdot  & \cdot  \\
        \end{pmatrix}
    $$
\end{frame}

\begin{frame}{MM Is Non-Commutative}
    MM is undefined if the inner dimensions differ.

    But even when the MM is defined, changing the order of the factors will change the result.

    \begin{boxed}
        For two square matrices $\m{A} \in \mathbb{R}^n, \m{B} \in \mathbb{R}^n,$ and $n > 1$, in general
        $$\m{A}\,\m{B} \ne \m{B}\,\m{A}$$
    \end{boxed}

    Example:
    $$\begin{pmatrix} 0 & 1 \\ 0 & 0 \end{pmatrix}\begin{pmatrix} 0 & 0 \\ 1 & 0 \end{pmatrix}=\begin{pmatrix} 1 & 0 \\ 0 & 0 \end{pmatrix}
        \quad \neq \quad
        \begin{pmatrix} 0 & 0 \\ 1 & 0 \end{pmatrix}\begin{pmatrix} 0 & 1 \\ 0 & 0 \end{pmatrix} = \begin{pmatrix} 0 & 0 \\ 0 & 1 \end{pmatrix}$$
\end{frame}

\begin{frame}{MM Is Distributive}
    \textbf{Distributivity}: MM is distributive with respect to matrix addition. If $\m{A}$, $\m{B}$, $\m{C}$, $\m{D}$ are
    matrices of size $m \times n$, $n \times p$, $n \times p$, , $p \times q$, respectively, we have
    \begin{itemize}
        \item left distributivity: $\m{A}(\m{B} + \m{C}) = \m{AB} + \m{AC}$
        \item right distributivity: $(\m{B} + \m{C} )\m{D} = \m{BD} + \m{CD}$
    \end{itemize}

    \textbf{Scalar multiplication}:
    \begin{itemize}
        \item For any $c \in \mathbb{R}$ we have $c\m{A} = \m{A}c$
        \item Also, $c(\m{AB}) = (c \m{A})\m{B}$ and $(\m{A} \m{B})c=\m{A}(\m{B}c)$
    \end{itemize}
\end{frame}


\begin{frame}{Matrix Transpose}
    \begin{boxed}
        \textbf{Transposition}\\
        For any matrix $\m{A} = [a_{ij}]$, its transpose is defined to be
        $$\m{A}^\top  = [a_{ji}]$$
    \end{boxed}

    Example:
    $$
        \begin{pmatrix}
            1 & 2 \\
            3 & 4 \\
            5 & 6
        \end{pmatrix}^{\operatorname{T}}
        =
        \begin{pmatrix}
            1 & 3 & 5 \\
            2 & 4 & 6
        \end{pmatrix}
    $$
\end{frame}

\begin{frame}{Matrix Transpose: Identities}
    Some useful matrix identities related to its transpose:
    \begin{boxed}
        \begin{itemize}
            \item $(\m{A^\top})^\top = \m{A}$
            \item $\left(\m{A} + \m{B}\right)^\top = \m{A}^\top + \m{B}^\top$
            \item $\left(\m{A}_1 \m{A}_2 \ldots \m{A}_{k-1} \m{A}_k\right)^\top = \m{A}_k^\top \m{A}_{k-1}^\top \ldots \m{A}_2^\top \m{A}_1^\top \quad \text{ e.g., }(\m{AB})^\top = \m{B}^\top\m{A}^\top$
            \item $(c\m{A})^\top = c \m{A}^\top$
            \item $\det \m{A} = \det \m{A}^\top$
            \item $\text{If}\, \forall i,j: a_{ij} \in \mathbb{R}\, \Rightarrow \,\m{A}^\top \m{A} \succeq 0 \quad \text{i.e., $\m{A}$ PSD}$
            \item $(\m{A}^{-1})^\top = (\m{A}^\top)^{-1}$
            \item $\text{If }\,\m{A} \in \mathbb{R}^{n\times n}\,\text{then eigenvalues unchanged under transposition}$
        \end{itemize}
    \end{boxed}
\end{frame}

\begin{frame}{Matrix Rank}
    Various definitions exists, all equivalent.
    \begin{boxed}
        We define
        $$\rank \m A := \text{number of linearly independent rows (or columns)}$$
    \end{boxed}

    Example: The matrix
    $$\begin{pmatrix}1&0&1\\-2&-3&1\\3&3&0\end{pmatrix}$$
    is of $\rank 2$ (do you see why?)
\end{frame}

\subsection{Determinants, Definiteness, Trace, Inverse}

\begin{frame}{Determinants}
    The determinant of a \emph{square} matrix $\m A$ is scalar that characterizes certain properties of $\m A$.
    Most importantly, the determinant is nonzero iff the matrix is invertible\\[5mm]

    Due to being expensive to compute, explicit use of determinants in machine learning applications is rare\\[5mm]

    Computation of the determinant of $\m A \in \mathbb{R}^n$:
    \begin{itemize}
        \item $n = 1$: $\det a = a$
        \item $n = 2$: $\det \begin{pmatrix} a & b\\c & d \end{pmatrix}=ad-bc$
        \item $n = 3$: $\det \begin{pmatrix} a & b & c \\ d & e & f \\ g & h & i \end{pmatrix}= aei + bfg + cdh - ceg - bdi - afh$
        \item $n > 3$: Usually computed via decomposition (see below)
    \end{itemize}
\end{frame}

\begin{frame}{Definiteness}
    For a \emph{symmetric} matrix $\m A \in \mathbb{R}^{n \times n}$ and
    $\v x \in \mathbb{R}^n \setminus \{\v 0\}$ consider $$s := \v x^\top \m A \v x.$$

    Then, $\m A$ is called
    \begin{itemize}
        \item \emph{positive-definite} (PSD) if $s > 0$ $\forall \v x$
        \item \emph{positive-semidefinite} (PD) if $s \geq 0$ $\forall \v x$
        \item \emph{negative-semidefinite} (NSD) if $s \leq 0$ $\forall \v x$
        \item \emph{negative-definite} (ND) if $s < 0$ $\forall \v x$
        \item and \emph{indefinite} otherwise
    \end{itemize}

    Example:
    The identity matrix $\m I = \begin{pmatrix} 1 & 0 \\ 0 & 1 \end{pmatrix}$ is positive-definite. To see this, let $\v x = \mqty(a \\ b)$ and consider
    $$\v{x}^\top \m I \v{x} = \begin{pmatrix} a & b \end{pmatrix}^\top \begin{pmatrix} 1 & 0 \\ 0 & 1 \end{pmatrix} \begin{pmatrix} a \\ b \end{pmatrix} = a^2 + b^2$$
    which is strictly positive for all $\v x \neq \v 0$.
\end{frame}

\begin{frame}{Matrix Trace}
    The trace of a \emph{square} matrix $\m A \in \mathbb{R}^{n \times n}$ is commonly defined as the sum of its diagonal elements:
    $$\tr \m{A} = \sum_{i=1}^n a_{ii} = a_{11} + a_{22} + \dots + a_{nn}.$$

    Example:
    $$
        \tr \begin{pmatrix}
            1  & 0  & 3  \\
            11 & 5  & 2  \\
            6  & 12 & -5
        \end{pmatrix} = 1 + 5 + (-5) = 1
    $$

    Properties:
    \begin{itemize}
        \item Linearity: $\tr(\m{A} + \m{B}) = \tr(\m{A}) + \tr(\m{B})$ and $\tr(c\m{A}) = c \tr(\m{A})$
        \item $\tr \m A = \tr \m A^\top$
        \item $\tr(\m{A}\m{B}\m{C}\m{D}) = \tr(\m{B}\m{C}\m{D}\m{A}) = \tr(\m{C}\m{D}\m{A}\m{B}) = \tr(\m{D}\m{A}\m{B}\m{C})$
              and as a special case $\tr(\m A \m B) =  \tr(\m B \m A)$
        \item Trace as the sum of eigenvalues $\lambda_i$ of $\m A$: $\tr(\mathbf{A}) = \sum_{i=1}^n \lambda_i$
    \end{itemize}

\end{frame}

\conceptcheck{Determine whether or not the real symmetric matrix $$\m A = \mqty(2 & -1 & 0 \\ -1 & 2 & -1 \\ 0 & -1 & 2)$$
    is positive-definite.}
{   $\m A$ is PD since for any non-zero column vector $\v z = \mqty(a & b & c)^\top$:
    $$
        \v{z}^\top \m A \v{z} = \left(\v{z}^\top \m M\right) \v{z}
        = \begin{pmatrix} (2a - b) & (-a + 2b - c) & (-b + 2c) \end{pmatrix}
        \begin{pmatrix} a \\ b \\ c \end{pmatrix} = \ldots =  a^2 + (a - b)^2 + (b - c)^2 + c^2 > 0
    $$
}

\subsection{Matrix Inversion and Decompositions}

\begin{frame}{Matrix Inverse}
    The inverse matrix $\m A^{-1}$ to a regular matrix $\m A$ is defined s.t.
    $$\m A \m A^{-1} = \m A^{-1} \m A = \m I$$
    with $\m I$ then identity matrix

    Example: The inverse of the real matrix
    $$\displaystyle \m A = \begin{pmatrix} 2 & 1 \\ 6 & 4 \end{pmatrix} \quad \text{is} \quad \m A^{-1} = \begin{pmatrix} 2 & -\tfrac{1}{2} \\ -3 & 1 \end{pmatrix}$$
    Verify this:
    $$
        \m A \m A^{-1} = \begin{pmatrix} 2 & 1 \\ 6 & 4 \end{pmatrix} \begin{pmatrix} 2 & -\tfrac{1}{2} \\ -3 & 1 \end{pmatrix}
        = \begin{pmatrix} 4-3 & -1+1 \\ 12-12 & -3+4 \end{pmatrix} = \begin{pmatrix} 1 & 0 \\ 0 & 1 \end{pmatrix} = \m I
    $$
\end{frame}

\begin{frame}{Matrix Inverse}
    When is a matrix $\m A$ invertible? Many equivalent conditions exist (``Invertible Matrix Theorem''):
    \begin{itemize}
        \item $\m A$ is row- or column-equivalent to the $n \times n$ identity matrix $\m{I}_n$ $\Rightarrow$ Gaussian Elimination
        \item The columns of A are linearly independent
        \item $\m A$ has full rank, $\rank \m A = n$
        \item $0$ is \emph{not} an eigenvalue of $\m A$ $\Rightarrow$ Definiteness
        \item $\det A \neq 0$
    \end{itemize}

    Inversion plays an important part in the solution of linear systems of equations:
    $$\m A \v x = \v b \quad \Longleftrightarrow \quad \v x = \m A^{-1} \v b$$
    Closely related to rank: For instance, if $\rank \m A = n$ (full rank), the system
    \begin{itemize}
        \item $\m A \v x = \v 0$ has only the trivial solution $\v x = 0$
        \item $\m A \v x = \v b$ has exactly one solution for each $\v b \in \mathbb{R}^n$
    \end{itemize}
    In practice, computing the inverse is costly and numerically delicate---avoid whenever possible
\end{frame}

\begin{frame}{Linear Systems}
    \begin{center}
        \begin{minipage}[b]{1.8in}
            \raisebox{-2pt}[8pt][0pt]{\small \begin{tabular}{@{}l}
                    \small \textit{Unique solution}
                \end{tabular}}
            \begin{center}
                \includegraphics{gr/mp/ch1.3} \\[.75ex]
                \small $\begin{linsys}{2}
                        3x  &+  &2y  &=  &7   \\
                        x   &-  &y   &=  &-1
                    \end{linsys}$
            \end{center}
        \end{minipage}
        \hspace*{0em}
        \begin{minipage}[b]{1.8in}
            \raisebox{-2pt}[8pt][0pt]{\small \begin{tabular}{@{}l}
                    \small \textit{No solutions}
                \end{tabular}}
            \begin{center}
                \includegraphics{gr/mp/ch1.4} \\[.75ex]
                \small $\begin{linsys}{2}
                        3x  &+  &2y  &=  &7   \\
                        3x  &+  &2y  &=  &4
                    \end{linsys}$
            \end{center}
        \end{minipage}
        \hspace*{0em}
        \begin{minipage}[b]{1.8in}
            \raisebox{-2pt}[8pt][0pt]{\small \begin{tabular}[t]{@{}l}
                    \textit{Infinitely many} \\
                    \textit{solutions}
                \end{tabular}}
            \begin{center}
                \includegraphics{gr/mp/ch1.5}         \\[.75ex]
                \small $ \begin{linsys}{2}
                        3x  &+  &2y  &=  &7   \\
                        6x  &+  &4y  &=  &14
                    \end{linsys}$
            \end{center}
        \end{minipage}
    \end{center}
\end{frame}

\conceptcheck{
    Find the coefficients $a$, $b$, and $c$ so that the graph of $f(x) = a x^2 + b x + c$ passes
    through the points $(1, 2)$, $(-1, 6)$, and $(2, 3)$
}{
    % \begin{verbatim}
    %     A = np.array([[1, 1, 1], [1, -1, 1], [4, 2, 1]])
    %     b = np.array([2, 6, 3])
    %     print(np.linalg.solve(A, b))  # => [ 1. -2.  3.]
    % \end{verbatim}
    $a=1$, $b=-2$, $c=3$
}

\begin{frame}{Matrix Inverse: Identities}
    Some useful matrix identities to inversion:
    \begin{boxed}
        \begin{itemize}
            \item $(\mathbf A^{-1})^{-1} = \mathbf A$
            \item $(k \mathbf A)^{-1} = k^{-1} \mathbf A^{-1}$ for any scalar $k \neq 0$
            \item $(\mathbf A^\top)^{-1} = (\mathbf A^{-1})^\top$
            \item For invertible matrices $\m A_i \in \mathbb{R}^{n \times n}$, $(\mathbf A_1 \mathbf A_2 \cdots \mathbf A_{k-1} \mathbf A_k)^{-1} = \mathbf A_k^{-1} \mathbf A_{k-1}^{-1} \cdots \mathbf A_2^{-1} \mathbf A_1^{-1}$
            \item $\det \mathbf A^{-1} = (\det \mathbf A)^{-1}$
        \end{itemize}
    \end{boxed}
\end{frame}

\begin{frame}{Matrix Decompositions}
    \framesubtitle{Eigendecomposition}
    The Eigendecomposition factorizes a diagonalizable matrix into a its canonical form,
    represented by \emph{eigenvalues} and \emph{eigenvectors}. Particularly important special case
    in machine learning: real, symmetric matrices

    \begin{boxed}
        A vector $\v x \in \mathbb{R}^n \setminus \{\v 0\}$ is an eigenvector of a square $N \times N$ matrix $\m A$, if it satisfies
        $$\m A \v x = \lambda \v x$$
        for some scalar $\lambda$. Then $\lambda$ is called the eigenvalue corresponding to eigenvector $\v x$.
    \end{boxed}

    \textbf{Computation}
    \begin{itemize}
        \item Characteristic polynomial $p(\lambda) := \det (\m A \m I - \lambda) = 0$ yields up to $N$ distinct eigenvalues $\lambda_i$
        \item Each $\lambda_i$ gives rise to a specific eigenvalue equation $(\m A  - \lambda \m I)\v v_i = \v 0$,
              yielding the eigenvectors $\v v_i$
    \end{itemize}

\end{frame}

\begin{frame}{Matrix Decompositions}
    \framesubtitle{Lower-upper (LU) Decomposition}

    \begin{boxed}
        The LU decomposition factors a square matrix $\m A$ into a lower triangular and an upper triangular matrix:
        $$\m A = \m L \m U$$
    \end{boxed}

    Example: An invertible $3 \times 3$ matrix $\m A$ can be factored as
    $$
        \begin{pmatrix}
            a_{11} & a_{12} & a_{13} \\
            a_{21} & a_{22} & a_{23} \\
            a_{31} & a_{32} & a_{33}
        \end{pmatrix} =
        \begin{pmatrix}
            \ell_{11} & 0         & 0         \\
            \ell_{21} & \ell_{22} & 0         \\
            \ell_{31} & \ell_{32} & \ell_{33}
        \end{pmatrix}
        \begin{pmatrix}
            u_{11} & u_{12} & u_{13} \\
            0      & u_{22} & u_{23} \\
            0      & 0      & u_{33}
        \end{pmatrix}.
    $$
    Remarks:
    \begin{itemize}
        \item Typical strategy used by computers to solve square systems of equations
        \item Often ``LU'' actually refers to  ``LUP'' where $\m P$ is a permutation matrix (``partial pivoting'')
        \item Any square matrix $\m A$ admits LUP and PLU factorizations. If $\m A$ is invertible, then it admits LU factorization
    \end{itemize}
\end{frame}

\begin{frame}{Matrix Decompositions}
    \framesubtitle{Lower-upper (LU) Decomposition}
    LU decompositions find wide-spread application from solving linear equations, to inverting a matrix,
    to computation of the determinant. \\[5mm]

    \textbf{Example}: Solve a system of linear equations with LUP\\
    Given the system $\m A \v x = \v b$ and the decomposition $\m P \m A = \m L \m U$, rewrite the problem
    $$\m A \v x = \v b \quad \Longleftrightarrow \quad \m P \m A \v x = \m P \v b \quad \Longleftrightarrow \quad \m L \m U \v x = \m P \v b$$
    Then $\v x$ can be found by
    \begin{enumerate}
        \item first solving $\m L \v y = \m P \v b$ for $\v y$ (forward substitution)
        \item second solving $\m U \v x = \v y$ for $\v x$ (backward substitution)
    \end{enumerate}
\end{frame}

\begin{frame}{Matrix Decompositions}
    \framesubtitle{Cholesky Decomposition}
    Why Cholesky? As a special case of LU for PD matrices, about twice as efficient to compute

    \begin{boxed}
        The Cholesky decomposition of an invertible matrix $\m A$ is defined as $$\m A = \m L \m L^\top$$
        where $\m L$ is a real lower triangular matrix with positive diagonal entries
    \end{boxed}

    Wide-spread applications:
    \begin{itemize}
        \item Linear least squares: efficient solutions of linear system arising from partial differential equations
        \item Non-linear optimization: numerical stability in Quasi-Newton by avoiding direct updates of the Hessian
        \item Monte Carlo simulations: produce correlated samples from uniform ones
        \item Kalman filters: forcing the system covariance matrix to remain positive semi-definite
    \end{itemize}
\end{frame}


\begin{frame}{Matrix Decompositions}
    \framesubtitle{Singular Value Decomposition (SVD)}
    SVD generalizes the eigendecomposition of a square matrix to any $m \times n$ matrix.

    \begin{boxed}
        The SVD decomposes a real $m \times n$ matrix $\m A$ as follows,
        $$\m A = \m U \m \Sigma \m V^\top$$
        with $\m U$ and $\m V$ are real orthogonal matrices
        and $\m \Sigma $ a $m \times n$ rectangular diagonal matrix w
    \end{boxed}

    If M is real, then U and V can be guaranteed to be real orthogonal matrices; in such contexts, the SVD is often denoted   $\displaystyle \mathbf {U\Sigma V} ^{\mathsf {T}}$
\end{frame}

\include{pt1_ex}
\section{Part 2: Calculus}

\subsection{Derivatives: Polynomials, Power Functions, Product and Chain Rule, Multivariate Calculus}

\begin{frame}{Intuition}
    \textbf{Geometrically}: The derivative is the tangent line to the curve at the point $(x0, y0)$
    is a line passing through $(x0, y0)$ and `flat against' the curve

    \begin{center}
        \includesvg[width=0.3\textwidth]{fig/tangent_to_a_curve}
    \end{center}

    \textbf{Intuition}: The derivative measures the sensitivity to change of the function value
    (output value) with respect to a change in its argument (input value)\\[5mm]

\end{frame}

\begin{frame}{Definition}
    Roughly speaking, a function $f(x)$ is differentiable at a point $a$ of its domain,
    if the limit
    $$L=\lim_{h \to 0}\frac{f(a+h)-f(a)}h$$
    exists.

    \begin{boxed}
        If the limit $L$ exists, then this limit is called the \emph{derivative} of $f$ at $a$,
        and denoted $$f'(a) \quad \text{or} \quad \frac{df}{dx}(a)$$
    \end{boxed}
\end{frame}



\begin{frame}{Derivatives of polynomials}
    There's just four simple facts which suffice to take the derivative
    of any polynomial plussomewhat more general things.

    First, there is the rule for taking the derivative of a \emph{power function}
    $f(x)=x^n$:
    $${d\over dx}x^n=n\,x^{n-1}$$

    Second, a special case of the rule above, the derivative of any constant $c$ is zero:
    $${d\over dx}\,c=0$$

    Third, for any function $f(x)$, multiplicative constants do not affect the result,
    $${d\over dx}\,c f=c \,{d\over dx}\,f$$

    Fourth, for any two functions $f(x)$ and  $g(x)$, the derivative of the sum
    is the sum of the derivatives:
    $${d\over dx}(f+g)={d\over dx}f+{d\over dx}g$$
\end{frame}

\begin{frame}
    Putting these four things together, we can write general formulas like
    $${d\over dx}(ax^m+bx^n+cx^p)=a\,m\,x^{m-1}+b\,n\,x^{n-1}+c\,p\,x^{p-1}$$
    and so on, with any number of summands.

    As a refresher, here are some examples:

    $${d\over dx}5x^3=15x^2$$
    $${d\over dx}(3x^7+5x^3-11)=21x^6+15x^2$$
    $${d\over dx}(2-3x^2-2x^3)=-6x-6x^2$$
    $${d\over dx}(-x^4+2x^5+ 1)=-4x^3+10x^4$$
\end{frame}

\begin{frame}{More general power functions}
    Remember some of the other possibilities for the exponential notation $x^n$. For example
    $$x^{{1\over 2}}=\sqrt{x}$$
    $$x^{-1}={1\over x}$$
    $$x^{-{1\over 2}}={ 1 \over \sqrt{x }}$$
    and so on.

    Good news: Above rule for taking
    the derivative of powers of $x$ is still correct here, even for
    exponents which are negative or fractions or even real numbers:

    $$ {d\over dx}x^r=r\,x^{r-1}$$
\end{frame}

\begin{frame}
    Thus, in particular,
    $${d\over dx}\sqrt{x}=
        {d\over dx}x^{{1\over 2}}={1\over 2}x^{-{1\over 2}}$$
    $${d\over dx}{1\over x}={d\over dx}x^{-1}=-1\cdot
        x^{-2}={-1\over x^2}$$

    When combined with the sum rule from above, we have the
    obvious possibilities:

    $${d\over dx}(3x^2-7\sqrt{x}+{5\over x^2})=
        {d\over dx}(3x^2-7x^{{1\over 2}}+5x^{-2})=
        6x-{7\over 2}x^{-{1\over 2}}-10x^{-3}$$

    \begin{boxed}
        Expressing square roots, cube roots, inverses,
        etc., in terms of exponents often simplifies
        the rearranging of algebraic expressions
    \end{boxed}

\end{frame}

\begin{frame}{Quotient rule}
    \begin{boxed}
        \textbf{Quotient rule}:
        Given two differentiable functions $f(x)$ and $g(x)$, the derivative of the
        ratio of the two is given by:
        $${d\over dx}{f\over g}={f'g-g'f\over g^2}$$
    \end{boxed}

    Example: For $f(x) = x-1$ and $g(x) = x-2$ we get
    $${d\over dx}{x-1\over x-2}=
        {(x-1)'(x-2)-(x-1)(x-2)'\over (x-2)^2}=
        {1\cdot(x-2)-(x-1)\cdot 1\over (x-2)^2}$$
    $$=
        {(x-2)-(x-1)\over (x-2)^2}=
        {-1\over (x-2)^2}$$
\end{frame}

\begin{frame}
    Special case of the quotient rule: Applying the quotient rule to the
    expression $1 / g(x)$ yields
    \begin{boxed}
        \textbf{Reciprocal rule}:
        $$\frac{d}{dx}\frac{1}{g(x)}=-\frac{g'(x)}{g(x)^2}$$
    \end{boxed}

    Because
    $$\frac{d}{dx}\frac{1}{g(x)}=\frac{0 \cdot g(x) - 1 \cdot g'(x)}{g(x)^2}=\frac{-g'(x)}{g(x)^2}$$
\end{frame}

\begin{frame}{Product rule}
    We have:
    \begin{boxed}
        \textbf{Product rule}:
        $${d\over dx}(fg)=f'g+fg'$$
    \end{boxed}
    Note the this rule behaves differently from the simple rule for sums above

    Example: Differentiate $f(x) = x^2 \sin x$. We get $f'(x) = 2 x \sin x + x^2 \cos x$

    Note: ``constant multiple rule'' is a special case: $(c f(x))' = c f'(x)$ since ${d \over dx} c = 0$
\end{frame}

\conceptcheck{Find the derivative of the sigmoid function
    $$\sigma(x) = \frac{1}{1 + e^{-x}}$$
    and express it using the sigmoid function itself.
    Hint: The identity $\frac{a}{1 + a} = \frac{a + 1 - 1}{1 + a} = 1 - \frac{1}{1 + a}$ should prove useful
}{
    $$\frac{d}{\d \sigma} \frac{1}{1 + e^{-x}} = \ldots
        = \frac{1}{1 + e^{-x}} \frac{e^{-x}}{1 + e^{-x}}
        = \frac{1}{1 + e^{-x}} \frac{1 + e^{-x} - 1}{1 + e^{-x} }
        = \frac{1}{1 + e^{-x}}\left( 1 - \frac{1}{1 + e^{-x} }\right)
        = \sigma(x) \, (1 - \sigma(x))
    $$
}

\begin{frame}{Chain rule}
    Function composition operator ``$\circ$'':\\
    Operator that takes two functions $f$ and $g$ and produces a function $h = g \circ f$
    such that $h(x) = g(f(x))$. In other words, $g$ is applied to the result of $f(x)$.\\[5mm]

    The chain rule states how the derivative of the composition of two
    differentiable functions $f$ and $g$ behaves:
    \begin{boxed}
        \textbf{Chain rule}:
        $$h'=(f\circ g)'=(f'\circ g)\cdot g'$$
        \center{or less formally}
        $${d\over dx} f(g(x))=f'(g(x))\cdot g'(x)$$
    \end{boxed}

    Note that the chain rule is a core utility in machine learning, e.g. forming the base of
    how neural networks are learning, among countless other uses.
\end{frame}

\begin{frame}
    Use of chain rule is often straightforward. Just be clear what $f(x)$ and what $g(x)$ are.

    Here are some examples.

    Example 1:
    $${d\over dx}(1+x^2)^{100}=f'(g(x))\cdot g'(x)=100\,g^{99}(x)\cdot 2x=
        100\,(1+x^2)^{99}\cdot 2x$$

    Example 2:
    $${d\over dx}\sqrt{3x+2}={d\over dx}(3x+2)^{1/2}=
            {1\over 2}(3x+2)^{-1/2}\cdot 3$$

    Example 3:
    $${d\over dx}(3x^5-x+14)^{11}=11\,(3x^5-x+14)^{10}\cdot(15x^4-1)$$

\end{frame}


\begin{frame}
    Similarly the chain rule applies to composites of more than two functions:
    $$
        (f \circ g \circ h)'
        = f' \circ (g \circ h)\cdot (g \circ h)'
        = f' \circ (g \circ h) \cdot g'(h) \cdot h'
        = (f' \circ g \circ h) \cdot (g' \circ h) \cdot h'
    $$
    \\[5mm]
    For example, compute ${d \over dx} e^{\sin x^2}$:
    \begin{itemize}
        \item Define individual functions and compute their derivatives:
              \begin{align*}
                  y := f(u) = e^u \quad               & \Longrightarrow \quad \frac{dy}{du} = f'(u) = e^u = e^{\sin x^2} \\
                  u := g(v) = \sin v = \sin x^2 \quad & \Longrightarrow \quad \frac{du}{dv} = g'(v) = \cos v = \cos x^2  \\
                  v := h(x) = x^2 \quad               & \Longrightarrow \quad \frac{dv}{dx} = h'(x) = 2x
              \end{align*}
        \item Therefore:
              $${d \over dx} e^{\sin x^2} = e^{\sin x^2} \cos(2x)\;2x$$
    \end{itemize}
\end{frame}

\begin{frame}{Multivariate Chain Rule}
    \begin{boxed}
        \textbf{Multivariate Chain Rule}:
        Let
        \begin{itemize}
            \item $z=f(x_1,x_2,\ldots,x_m)$ be differentiable in $m$ independent variables, and let
            \item $\forall i \in 1, \ldots,m: x_i=x_i(t_1,t_2,\dots,t_n)$ be differentiable in
                  $n$ independent variables
        \end{itemize}

        Then, $\forall j \in 1,2,\ldots,n$:

        $${\partial z \over \partial t_j} =
            \sum_{i = 1}^m {\partial z \over \partial x_i } { \partial x_i \over \partial t_j }
            =
            { \partial z \over \partial x_1 } {\partial x_1 \over \partial t_j}
            + { \partial z \over \partial x_2 } {\partial x_2 \over \partial t_j} + \cdots
            + { \partial z \over \partial x_m } {\partial x_m \over \partial t_j}
            = \nabla z \cdot {\partial \v x \over \partial t_j }
        $$
    \end{boxed}

    For example, in the case of $x(t)$, $y(t)$, and $z=f(x,y)$, we get
    $${dz \over dt} = {\partial z \over \partial x} {dx \over dt} + {\partial z \over \partial y} {dy \over dt}$$
    with $z=f(x(t),y(t))$ being differentiable at $t$.
\end{frame}

\begin{frame}
    \textbf{Example: Multivariate Chain Rule}

    Given $u(x, y) = x^2 + 2y$ with $x(r, t) = r \sin(t)$ and $y(r,t) = \sin^2(t)$, we are interested in
    finding $\partial u / \partial r$ and $\partial u / \partial t$.

    $$\frac{\partial u}{\partial r}
        =\frac{\partial u}{\partial x} \frac{\partial x}{\partial r}+\frac{\partial u}{\partial y} \frac{\partial y}{\partial r}
        =2x \sin t + 2 \times 0 = 2 r \sin^2 t$$
    and
    \begin{align*}
        \frac{\partial u}{\partial t} & = \frac{\partial u}{\partial x} \frac{\partial x}{\partial t}+\frac{\partial u}{\partial y} \frac{\partial y}{\partial t} \\
                                      & = 2x \times r\cos t + 2 \times  2 \sin t \cos t                                                                           \\
                                      & = 2r \sin t \times  r \cos t + 4 \sin t \cos t                                                                            \\
                                      & = 2(r^2 + 2) \sin t \cos t                                                                                                \\
                                      & = (r^2 + 2) \sin(2t)
    \end{align*}
\end{frame}

\conceptcheck{You are given 200 m of wire to make a fence around some rectangular area of side length $a$
    and $b$, respectively. How should you choose $a$ and $b$ s.t. the area which you are covering is maximized?}{
    $$\max_{a,b} a \, b \quad \text{s.t.} \quad 2a+2b=200
        \quad \Longleftrightarrow \quad \max_a a\, (100 - a)
        \quad \Longleftrightarrow \quad a^\star = 50\quad \text{(square area)}$$
}

\subsection{Mathematical Optimization and Newton's Method}

\begin{frame}{Extreme Value Problems}
    Systematic procedure to find the minimum and maximum values of a
    function $f(x)$ on an interval $[a,b]$:

    \begin{enumerate}
        \item Solve $f'(x)=0$ to find the list of critical points of $f$
        \item Exclude any critical points not inside the interval $[a,b]$
        \item Add to the list the {\it endpoints} $a,b$ of the interval and
              any points of discontinuity or non-differentiability
        \item At each point on the list, evaluate $f$: the
              biggest number that occurs is the maximum, and the littlest number
              that occurs is the minimum.
    \end{enumerate}
    \vspace*{5mm}

    Example: Find the minima and maxima of the
    function $f(x)=x^4-8x^2+5$ on the interval $[-1,3]$. \\
    \begin{enumerate}
        \item $f'(x) = 4 x^3-16x = 0 \quad \Longleftrightarrow \quad x(x^2 - 4) = 0 \quad \Longleftrightarrow \quad x(x + 2)(x - 2)=0$\\
              So critical points are $-2,0,2$
        \item Remove $-2$ from list since outside interval
        \item Add endpoints $-1,3$ to list. Now list is $-1,0,2,3$
        \item Evaluate $f$ at $-1,0,2,3$ which yields $-2, 5, -11, 14$, respectively
    \end{enumerate}
    Hence maximum is $14$, which occurs at $x=3$, and the minimum is $-11$,
    which occurs at $x=2$.
\end{frame}


\begin{frame}{Global and Local Extrema}
    \begin{columns}[onlytextwidth]
        \begin{column}{0.6\textwidth}
            Use the \emph{derivative test} to determine the type of each cricitcal point.

            In the scalar case:
            \begin{itemize}
                \item $f''(x) < 0$: local maximum
                \item $f''(x) > 0$: local minimum
                \item $f''(x) = 0$: test inconclusive $\Rightarrow$ higher-order derivative test
            \end{itemize}

        \end{column}
        \begin{column}{0.4\textwidth}
            \includesvg[width=\textwidth]{fig/extrema_example}
        \end{column}
    \end{columns}

    Remarks:
    \begin{itemize}
        \item There is other tests to determine the nature of a critical point, e.g. first derivative test
        \item The above holds only for ``reasonably well-behaved'' functions
        \item The rules and intuition readily extend to the multi-dimensional case
    \end{itemize}


\end{frame}

\begin{frame}{Mathematical Optimization}
    \begin{boxed}
        \textbf{General optimization problem}:
        Given: a function $f : A \rightarrow \mathbb{R}$ from some set $A$ to the real numbers\\[1mm]
        Objective: find an element $\v x^\star \in A$ such that
        \begin{itemize}
            \item $\forall \v x \in A : f(\v x^\star) \leq f(\v x)$ \;\; ``minimization''
            \item $\forall \v x \in A : f(\v x^\star) \geq f(\v x)$ \;\; ``maximization''
        \end{itemize}
    \end{boxed}

    Optimization problems
    \begin{itemize}
        \item Two major classes: discrete and continuous optimization (focus on latter)
        \item We can always use \emph{minimization} because
              $$f(\mathbf{x^\star})\geq f(\mathbf{x}) \quad \Longleftrightarrow \quad -f(\mathbf{x^\star})\leq -f(\mathbf{x})$$
    \end{itemize}
\end{frame}

\begin{frame}{Optimization in Practice}
    In practice, algorithms that solve optimization problems fall into three broad classes:
    \begin{itemize}
        \item those that terminate in a finite number of steps, e.g. Simplex, combinatorial algos
        \item iterative methods that converge to a solution, e.g. Gradient descent, Newton's method, SQP
        \item heuristics that provide approx. solutions, e.g. genetic and evolutionary algos, PSO
    \end{itemize}

    Relevance for machine learning:
    \begin{itemize}
        \item Many machine learning problems are formulated as the
              minimization of some (so-called) loss function over a training set of examples
        \item Finding a solution may require to employ method from any of the three classes above
    \end{itemize}
\end{frame}

\begin{frame}
    \begin{boxed}
        \small
        \textbf{Nabla Operator}:
        Gradients in $\mathbb{R}^n$ are conveniently expressed with the \emph{Nabla operator}:
        $$
            {\nabla \equiv \mqty(\partial / \partial x_1 \\ \partial / \partial x_2 \\ \vdots \\ \partial / \partial x_n)  }
            \quad \text{or applied to some $f(\v x)$} \quad
            {\nabla f(\v x) = \mqty(\partial f/ \partial x_1 \\ \partial f / \partial x_2 \\ \vdots \\ \partial f / \partial x_n)  }
        $$
    \end{boxed}
    \vspace*{-1mm}

    \begin{boxed}
        \small
        \textbf{Hessian Matrix}:
        The Hessian is a square matrix of 2nd order partial derivatives of a scalar-valued function and a measure of local curvature:
        $$
            H_f= \begin{pmatrix}
                \dfrac{\partial^2 f}{\partial x_1^2}             & \dfrac{\partial^2 f}{\partial x_1\,\partial x_2} & \cdots & \dfrac{\partial^2 f}{\partial x_1\,\partial x_n} \\[2.2ex]
                \dfrac{\partial^2 f}{\partial x_2\,\partial x_1} & \dfrac{\partial^2 f}{\partial x_2^2}             & \cdots & \dfrac{\partial^2 f}{\partial x_2\,\partial x_n} \\[2.2ex]
                \vdots                                           & \vdots                                           & \ddots & \vdots                                           \\[2.2ex]
                \dfrac{\partial^2 f}{\partial x_n\,\partial x_1} & \dfrac{\partial^2 f}{\partial x_n\,\partial x_2} & \cdots & \dfrac{\partial^2 f}{\partial x_n^2}
            \end{pmatrix}
        $$
    \end{boxed}

\end{frame}

\begin{frame}{Gradient Descent}
    \begin{columns}[onlytextwidth]
        \begin{column}{0.65\textwidth}
            Gradient Descent (GD) is a first-order iterative optimization algorithm for finding a
            local minimum of a differentiable function:
            \begin{itemize}
                \item Assume $F(\v x)$ is defined and locally differentiable around $\v a$
                \item Then $F(\v x)$ decreases fastest in the direction
                      of the \emph{negative} gradient of $F$ at $a$, which is $-\nabla F(\v a)$
                \item For a sufficiently small step size (learning rate) $\gamma \in \mathbb{R}^+$,
                      we have
                      $$
                          \v a_{n+1}=\v a_{n}-\gamma \nabla F(\v a_{n})
                          \quad \text{and} \quad F(\mathbf{a_{n}} )\geq F(\mathbf{a_{n+1}})
                      $$
            \end{itemize}
        \end{column}
        \begin{column}{0.35\textwidth}
            \begin{flushright}
                \includesvg[width=.7\textwidth]{fig/gradient_descent}
            \end{flushright}
        \end{column}
    \end{columns}

    Iterating the above we get the
    \begin{boxed}
        \textbf{Gradient Descent Update Rule}:
        $$
            \v x_{n+1}=\v x_n-\gamma_n \nabla F(\v x_n),\ n \ge 0
        $$
        where $F(\v x_0)\ge F(\v x_1)\ge F(\v x_2)\ge \ldots$ is a monotonic sequence
    \end{boxed}

\end{frame}

\begin{frame}
    Note: Convergence is \emph{guaranteed} under certain reasonable assumptions and
    a proper choice of $\gamma$ (e.g. via line search or the Barzilai Borwein method)\\[5mm]

    Wide range of applications: Linear and non-linear systems, least squares problems, anything ``minimizable''
\end{frame}

\begin{frame}{Newton's Method}
    \begin{columns}[onlytextwidth]
        \begin{column}{0.65\textwidth}
            Iterative prodecure for finding roots $\v x^\star$ of a differentiable function $F$:
            $$\{ \v x_i^\star : F(\v x_i^\star) = 0 \}$$
            Intuitively, does so by repeatedly intersecting tangents with $x$-Axis:
            $$y=F(x_k) + F'(x_k)(x-x_k) \overset{!}{=} 0 \quad \Longrightarrow \quad x_{k+1} = x_k - \frac{F(x_k)}{F'(x_k)}$$

        \end{column}
        \begin{column}{0.35\textwidth}
            \begin{flushright}
                \includesvg[width=.6\textwidth]{fig/newton_optimization_vs_grad_descent}
            \end{flushright}
        \end{column}
    \end{columns}

    Applicable to minimization problem $\min_{x \in \mathbb{R}} f(x)$:
    Simply let $F(x) = f'(x)$ which yields:
    \begin{boxed}
        \textbf{Newton's Method (for optimization)}:
        $$x_{k+1} = x_k - \frac{f'(x_k)}{f''(x_k)} \quad \text{(univariate case)}$$
    \end{boxed}

    Corresponds to a second-order approxmiation of $f$ $\Rightarrow$ fast convergence
\end{frame}

\begin{frame}
    \vspace*{5mm}
    Really Newton is just a second-order Taylor expansion of $f$ around $x_k$:

    $$f(x_k + t) \approx  f(x_k) + f'(x_k) t + \frac{1}{2} f''(x_k) t^2$$
    and minmizing its first derivative:
    $${d \over dt} f(x_k + t) \approx  f'(x_k) + f''(x_k) t \overset{!}{=} 0 \quad \Longrightarrow
        \quad x_{k+1} = x_k - \frac{f'(x_k)}{f''(x_k)}
    $$
    Similarly in the multi-variate case, with $\m H_f$ the Hessian of $f$:
    $$
        f(\v x_k + \v a) \approx f(\v x_k) + \nabla f(\v x_k)^\top \v a + \frac12 \v a^\top \m H \v a
        \quad \text{and} \quad
        {\partial \over \partial \v a} f(\v x_k + \v a) \approx \nabla f(\v x_k) + \m H \v a \overset{!}{=} \v 0
    $$

    \begin{boxed}
        \textbf{Newton's Method (for optimization)}:
        $$\v x_{k+1} = \v x_k - \m H_f^{-1}\,\nabla f(\v x_k) \quad \text{(multivariate case)}$$

        Note that this yields the exact solution for any quadratic objective function in a \emph{single} step.
    \end{boxed}
\end{frame}

\subsection{Stochastic Gradient Descent}

\begin{frame}{Stochastic Gradient Descent (SGD)}
    SGD is:
    \begin{itemize}
        \item an iterative optimization method, stochastic approximation of gradient descent
        \item the fundamental idea behind many machine learning algorithms such as linear
              support vector machines, logistic regression, graphical models
        \item the de facto standard training algorithm for neural networks
    \end{itemize}

    Given an objective function $Q_i(\v w)$ which computes some sort of error for sample $i$,
    e.g. $Q_i(\v w) = (\hat y_i - y_i)^2$

    \begin{boxed}
        \textbf{Classical SGD}:
        \input{algo/2_basic_sgd}
    \end{boxed}
\end{frame}

\begin{frame}
    \textbf{Example:} SGD to compute a least squares fit of a straight line\\
    Given training data ${(x_i, y_i)}_{i=1}^N$, fit the model
    $$\hat{y} = \! w_1 + w_2 x$$

    Objective function (minimize):
    $$Q(w) = \sum_{i=1}^n Q_i(w) = \sum_{i=1}^n \left(\hat{y_i}-y_i\right)^2 = \sum_{i=1}^n \left(w_1 + w_2 x_i - y_i\right)^2$$

    The update step of the Basic SGD algorithm above becomes
    $$\begin{bmatrix} w_1 \\ w_2 \end{bmatrix} \leftarrow
        \begin{bmatrix} w_1 \\ w_2 \end{bmatrix}
        - \eta \begin{bmatrix} \frac{\partial}{\partial w_1} (w_1 + w_2 x_i - y_i)^2 \\
            \frac{\partial}{\partial w_2} (w_1 + w_2 x_i - y_i)^2\end{bmatrix} =
        \begin{bmatrix} w_1 \\ w_2 \end{bmatrix}
        -  \eta  \begin{bmatrix} 2 (w_1 + w_2 x_i - y_i) \\ 2 x_i(w_1 + w_2 x_i - y_i) \end{bmatrix}.
    $$
\end{frame}

\begin{frame}{SGD Extensions}
    \textbf{Extension 1: Batch updates} \\
    Issue: Each iteration of classical SGD evaluates the gradient at a single point $x_i$ only\\
    $\quad \Rightarrow $ often results in noisy estimates of the true gradient\\
    $\quad \Rightarrow $ computationally expensive (blocking vectorization)

    Idea: Combine multiple gradient estimates of a so-called ``mini-batch'', then do a single update\\
    $\quad \Rightarrow $ smoother convergence thanks to better estimates \\
    $\quad \Rightarrow $ speed-up owing to vectorization
\end{frame}

\begin{frame}
    \vspace*{5mm}
    \textbf{Extension 2: Implicit updates (ISGD)}\\

    Issues:
    \begin{itemize}
        \item Classical SGD sensitive to choice of $\eta$ and may easily diverge
        \item Better convergence with some decay schedule of $\eta$ over iterations, $\eta = \eta(\text{iter})$
    \end{itemize}

    Idea: Evaluate stochastic gradient at the next iteration rather than the current one:
    $$\v w_{k + 1} = \v w_k - \eta \nabla Q_i(\v w_{k + 1})$$

    To see the effect, assuming an ordinary least squares objective
    $$Q_i(\v w) = \frac12 (\hat y_i - y_i)^2$$
    in the usual samples $\{(x_i, y_i)\}_{i=1}^N$, compare the resulting update rules:
    \begin{align*}
        \text{Classical SGD:}               & \quad \v w_{k + 1} = \v w_k + \eta (y_i - \v x_i^\top \v w_k) \v x_i \\
        \text{Implicit SGD (closed-form!):} & \quad \v w_{k + 1} =
        \v w_k + \frac{\eta}{1 + \eta \left\|\v x_i\right\|^2} \left(y_i - \v x_i^\top \v w_k\right) \v x_i
    \end{align*}

    Note the normalizing effect on $\eta$, resulting in numerically stable procedure for virtually all $\eta$
\end{frame}

\conceptcheck{
    You are in the middle of a minimization with ISGD:
    \begin{itemize}
        \item Iteration $k = 348$, number of samples $n = 50'000$
        \item Current weight vector: $\v w_{348} = \mqty(5 & -3)^\top$
        \item Current sample: $\v x_{1283} = \mqty(3 & -4)^\top, y_{1283} = 57$
        \item Learning rate: $\eta = 0.04$
    \end{itemize}

    Continue the iteration for one more step to find the next weight vector $\v w_{349}$.
}{
    \begin{align*}
        \v w_{349} & = \v w_{348} + \frac{\eta}{1 + \eta \left\|\v x_i\right\|^2} \left(y_i - \v x_i^\top \v w_k\right) \v x_i \\
                   & = \mqty(5 \\ -3) + \frac{0.04}{1 + 0.04 \times 25} \left(57 - 27 \right) \mqty(3 \\ -4) \\
                   & = \mqty(6.8 \\ -5.4)
    \end{align*}
}

\begin{frame}
    \vspace*{5mm}
    \textbf{Extension 3: Momentum}\\

    Issue: Classical SGD frequently shows oscillations, hampering convergence

    Idea: Introduce some ``memory'' $\Delta \v w_k$ of the current direction, i.e., do not take
    ``hard turns'' at every step:
    $$\Delta \v w_{k + 1} = -\eta \nabla Q_i(\v w_k) + \alpha \Delta \v w_k $$
    resulting in the update rule
    $$\v w_{k+1} =  \v w_k + \Delta \v w_{k + 1} = \v w_k - \eta \nabla Q_i(\v w_k) + \alpha \Delta \v w_k $$
    with $\alpha \in [0, 1)$ an exponential decay factor (hyperparameter)
\end{frame}

\begin{frame}
    \vspace*{5mm}
    \textbf{More Extensions}\\

    More sophisticated extensions that adapt the learning rate exists:
    \begin{itemize}
        \item AdaGrad (adaptive gradient): offers per-parameter learning rate
        \item Adam (adaptive momentum): also per-parameter learning rate; basis for most modern algorithms (AdaMax, AMSGrad)
        \item Second-order methods: stochastic analogue of the standard (deterministic) Newton algorithm
    \end{itemize}

\end{frame}

\subsection{Integral Calculus}

\begin{frame}{Integration}
    \begin{columns}[onlytextwidth]
        \begin{column}{0.73\textwidth}
            \textbf{Definite integrals} of real-valued functions $f(x)$ w.r.t. a real variable $x$ on an interval $[a, b]$ is written as
            $$\int_a^b f(x) \d x$$
            with limits $a$ and $b$. A function $f$ is \emph{integrable} if its integral over its domain is finite\\[5mm]

            \textbf{Indefinite integrals (antiderivative)}, written
            $$F(x) := \int f(x) \d x$$
            represent a class of functions whose derivative is the integrand, hence
            $$F' = f$$
        \end{column}
        \begin{column}{0.2\textwidth}
            \includegraphics[width=1\textwidth]{fig/riemann-2.eps}
            \includegraphics[width=1\textwidth]{fig/riemann-4.eps}
            \includegraphics[width=1\textwidth]{fig/riemann-7.eps}
        \end{column}
    \end{columns}
    \vfill
    {
        \tiny
        Please note that all statemets around \emph{integrals} and \emph{integrability} are meant in the \emph{Riemann} sense
    }
\end{frame}
\begin{frame}
    Property 1 (linearity): The integral of a linear combination is the linear combination of the integrals:
    $$\int_a^b (\alpha f + \beta g)(x) \d{x} = \alpha \int_a^b f(x) \d{x} + \beta \int_a^b g(x) \d{x}$$

    Property 2: Flipping integration bounds flips the sign of the result,
    $$\int_a^b f(x) \, \d x = - \int_b^a f(x) \, \d x$$

    Property 3: Given $\int_a^b f(x) \d x$ exists with $a \leq b$, then for any $c \in [a, b]$, we have
    $$\int_a^b f(x) \, \d x = \int_a^c f(x) \, \d x + \int_c^b f(x) \, \d x$$
\end{frame}

\begin{frame}{Closed-form Integration}
    Refresher by example:
    \begin{columns}[onlytextwidth]
        \begin{column}{0.5\textwidth}
            \begin{itemize}
                \item $\displaystyle{ \int a \d x = a x + C}$
                \item $\displaystyle{ \int x^n\,dx = \frac{x^{n+1}}{n+1} + C, \quad n\neq -1 }$
                \item $\displaystyle{ \int (ax + b)^n \, dx= \frac{(ax + b)^{n+1}}{a(n + 1)} + C, \quad n\neq -1}$
                \item $\displaystyle{ \int {1 \over x}\,dx = \ln \left|x \right| + C }$
                \item $\displaystyle{ \int e^{ax}\,dx = \frac{1}{a}e^{ax} + C }$
                \item $\displaystyle{ \int f'(x)e^{f(x)}\,dx = e^{f(x)} + C }$
                \item $\displaystyle{ \int a^x\,dx = \frac{a^x}{\ln a} + C }$
            \end{itemize}
        \end{column}
        \begin{column}{0.5\textwidth}
            \begin{itemize}
                \item $\displaystyle{ \int \ln x\,dx = x \ln x - x + C }$
                \item $\displaystyle{ \int \log_a x\,dx = \frac{x\ln x - x}{\ln a} + C }$
                \item $\displaystyle{ \int \sin{x}\, dx = -\cos{x} + C }$
                \item $\displaystyle{ \int \cos{x}\, dx = \sin{x} + C }$
                \item $\displaystyle{ \int \tan{x} \, dx = \ln{\left| \sec{x} \right|} + C = -\ln{\left| \cos {x} \right|} + C }$
                \item $\displaystyle{ \int_0^\infty \sqrt{x}\,e^{-x}\,dx = \frac{1}{2}\sqrt \pi } \quad \text{Gamma function}$
                \item $\displaystyle{ \int_0^\infty e^{-a x^2}\,dx = \frac{1}{2} \sqrt \frac {\pi} {a} \quad \text{Gaussian integral}, a>0}$
            \end{itemize}
        \end{column}
    \end{columns}
\end{frame}

\begin{frame}{Integration By Parts}
    If $y(x)=u(x)\,v(x)$, then
    $$y'=(uv)'=u'v+uv' \quad \Longleftrightarrow \quad uv'=(uv)'-u'v  $$

    Therefore we get:
    \begin{boxed}
        \textbf{Integration by parts} (helpful to integrate products of functions):
        $$\int u v' \d x = \int (uv)' \d x - \int u'v \d x = u v  - \int v u'  \d x$$
    \end{boxed}

    Example: To find the antiderivative $\int x\sin(x)dx$, we \emph{choose} $u := x$ and $v := \sin x$,
    $$\int x \sin x \d x = -x \cos x - \int(-\cos x ) \d x= \sin x - x \cos x +C$$

    Note that the ``right'' choice of $u$ and $v$ is crucial
\end{frame}

\begin{frame}{Integration By Substitution}
    Transforms one integral into another integral that is easier to compute

    \begin{boxed}
        \textbf{Integration by substitution} (helpful to integrate compositions of functions):\\
        To transform $\int f(g(x)) \d x$
        \begin{itemize}
            \item choose a \emph{suitable} substitution $u = g(x)$
            \item compute $\d u = g'(x) \d x$ \quad (not obvious)
            \item substitute integrand and differential, then solve the (simpler) integral
            \item back-substitute
        \end{itemize}
    \end{boxed}

    Example: To find the antiderivative of $x \cos(x^2+1)$, choose $u := x^2 + 1$, $\d u = 2 \d x$. Then
    $$
        \int x \cos(x^2+1) \d x
        = \frac{1}{2} \int 2x \cos(x^2+1) \d x
        = \frac{1}{2} \int\cos u \d u
        = \frac{1}{2}\sin u + C
        = \frac{1}{2}\sin(x^2+1) + C
    $$
\end{frame}

\begin{frame}{Numerical Integration}
    Issue: In practice, many integrals possess no closed-form solution

    Resort: Quadrature (numerical integration) for 1-dimensional integrals works well, e.g.
    $$\int\limits_a^b f(x) \d x \approx \sum_{i=1}^nf(x_i^*)\,\Delta x$$

    Multi-dimensional integrals (arising e.g. as normalizers in Bayesian stats) require more
    sophisticated techniques such as Monte Carlo integration. Popular algorithms in this class
    include Metropolis-Hastings and Gibbs sampling (both instances of
    Markov Chain Monte Carlo)
\end{frame}

% \begin{frame}
%     - special functions: Step, Sigmoid, Logit, ReLU \\
%     - lagrange multipliers \\
%     - practical constrained optimization, e.g. SQP \\
% \end{frame}

\include{pt2_ex}
% \begin{columns}[onlytextwidth]
%     \begin{column}{0.5\textwidth}
%     \end{column}
%     \begin{column}{0.5\textwidth}
%     \end{column}
% \end{columns}

\section{Part 3: Statistics and Probability}
\subsection{Random Variables}

{
    \small

    \begin{frame}{Introduction}
        Background
        \begin{itemize}
            \item Probability is the mathematical field concerned with reasoning under uncertainty
            \item Probabilistic models enable us to reason about the \emph{likelihood} of various events
            \item Use of probabilities to describe the frequencies of repeatable events
                  (like coin tosses) is fairly uncontroversial
        \end{itemize}

        Schools of thoughts
        \begin{itemize}
            \item Frequentists scholars adhere to an interpretation of probability that applies only to such repeatable events
            \item Bayesian scholars use the language of probability more broadly to formalize our reasoning under uncertainty
        \end{itemize}

        Bayesian probability is characterized by two unique features:
        \begin{itemize}
            \item assigning prior belief e.g., $P[\text{"moon is made out of cheese"}]$, opions may differ
            \item evidence, that is hard facts that we measure
        \end{itemize}
        Bayesian probability provides unambiguous rules for how we should \emph{update} our belief in light
        of new evidence, while allows for different individuals to start off with different prior beliefs
    \end{frame}

    \begin{frame}{Terminology}
        Imagine the experiment:
        \begin{itemize}
            \item Toss a coin $n$ times. The probability for heads is $p_h = 1/2$.
            \item We observe $n_h$ heads and $n_t = n - n_h$ tails
            \item Fraction of heads that we \emph{expect} to see should exactly match the expected fraction of tails
            \item The larger $n$ the less likely $n_h = n_t$ becomes
        \end{itemize}

        Terminology:
        \begin{itemize}
            \item $p_h$ is called a \emph{probability}, captures the certainty with which any given toss will come up heads
            \item Probabilities assign scores between $0$ and $1$ to \emph{events} (i.e. the outcomes of interest)
            \item Here the event of interest is ``heads'' and we denote the corresponding probability $P(\text{heads})$
            \item $P(\cdot) = 1$ indicates absolute certainty, and $P(\cdot) = 0$ impossibility
            \item $n_h / n$ and $n_t / n$ are called \emph{statistics} (not probabilities)---empirical quantities
                  computed based on the observed data
            \item \emph{Estimators} are special statistics that use the data to produce estimates of model parameters (like probabilities).
                  \emph{Consistent estimators} converge to the corresponding probability (cf. next slide)
        \end{itemize}
    \end{frame}

    \begin{frame}
        \vspace*{5mm}
        Estimators and Consistency
        \begin{center}
            \includegraphics[width=0.6\textwidth]{fig/prob_estimator_convergence.eps}
        \end{center}
    \end{frame}

    \begin{frame}{Formal Definition}
        Denote the set of possible outcomes $\mathcal{S}$, and call it the \emph{sample space} or \emph{outcome space}, e.g.:
        \begin{itemize}
            \item Rolling a single coin: $\mathcal{S} = \{\textrm{heads}, \textrm{tails}\}$
            \item Rolling a single die, $\mathcal{S} = \{1, 2, 3, 4, 5, 6\}$
            \item flipping two coins, possible outcomes are
                  $\{(\textrm{heads}, \textrm{heads}), (\textrm{heads}, \textrm{tails}), (\textrm{tails}, \textrm{heads}),  (\textrm{tails}, \textrm{tails})\}$
        \end{itemize}
        \emph{Events} are subsets of the sample space.

        \begin{boxed}
            \textbf{Axioms of Probability}:
            A \emph{probability function} maps events onto real values
            ${P: \mathcal{A} \subseteq \mathcal{S} \rightarrow [0,1]}$.
            The probability of an event $\mathcal{A}$ in the given sample space $\mathcal{S}$, denoted $P(\mathcal{A})$,
            satisfies the following properties (Kolmogorov, 1933):

            \begin{itemize}
                \item The probability of any event $\mathcal{A}$ is a non-negative real number, i.e., $P(\mathcal{A}) \geq 0$;
                \item The probability of the entire sample space is $1$, i.e., $P(\mathcal{S}) = 1$;
                \item For any countable sequence of events $\mathcal{A}_1, \mathcal{A}_2, \ldots$ that are \emph{mutually exclusive} ($\mathcal{A}_i \cap \mathcal{A}_j = \emptyset$ for all $i \neq j$), the probability that any of them happens is equal to the sum of their individual probabilities, i.e., $P(\bigcup_{i=1}^{\infty} \mathcal{A}_i) = \sum_{i=1}^{\infty} P(\mathcal{A}_i)$.
            \end{itemize}
        \end{boxed}
    \end{frame}



    \begin{frame}{Random Variables}
        A \emph{random variable} (RV) is a mapping from an underlying sample space to a set of values
        \begin{itemize}
            \item Random variables can be much coarser than the raw sample space. E.g. define RV
                  ``greater than 0.5'' is binary but has an infinite sample space
            \item Multiple RVs may share the same underlying sample space. E.g. for RVs ``home alarm goes off''
                  and ``my house was burglarized'', knowing the value taken by one random variable can tell us
                  something about the likely value of another random variable
        \end{itemize}

        Every value taken by a random variable corresponds
        to a subset of the underlying sample space.
        Thus the occurrence where the random variable $X$
        takes value $v$, denoted by $X=v$, is an \emph{event}
        and $P(X=v)$ denotes its probability.
    \end{frame}

    \begin{frame}
        Notation in practice
        \begin{itemize}
            \item $P(X)$ may refer broadly to the \emph{distribution} of $X$
            \item $P(X,Y) = P(X) P(Y)$, may be shorthand to express a statement
                  that is true for all of the values
                  that the random variables $X$ and $Y$, i.e.,
                  $\forall i,j : P(X=i \textrm{ and } Y=j) = P(X=i)P(Y=j)$
            \item $P(v)$ may be used when the random variable is clear from the context
            \item $P(1 \leq X \leq 3)$ may denote the probability of the event $\{1 \leq X \leq 3\}$
        \end{itemize}

        In general, there's two types of random variables:
        \begin{itemize}
            \item Discrete, e.g. flips of a coin or tosses of a die
            \item Continuous, e.g. height $H$ of a person. Note that in this case, asking e.g.
                  for $P(H = 1.8796549851915233234)$ does not make sense. Instead, we work with
                  probability \emph{densities} p
        \end{itemize}

    \end{frame}

    %   \begin{frame}{// additional}
    %       \url{http://d2l.ai/chapter_preliminaries/probability.html#multiple-random-variables}

    %       \url{http://d2l.ai/chapter_preliminaries/probability.html#expectations}
    %   \end{frame}

    \subsection{Density and Distribution Functions}

    \begin{frame}{Probability Density Functions (PDFs)}
        \begin{columns}[onlytextwidth]
            \begin{column}{0.55\textwidth}
                \begin{boxed}
                    \textbf{PDF}:
                    Informally, a PDF $p(x)$ is characterized by the following properties:
                    \begin{itemize}
                        \item ${\displaystyle p(x) \geq 0 }$
                        \item ${\displaystyle \int_{-\infty}^{\infty} p(x)  \d x = 1 }$
                    \end{itemize}
                    From this it follows that the probability of any single value $x$
                    is $0$, $P(X = x) = 0$ and that $p(x) > 1$ is admissible
                \end{boxed}

                To compute actual probabilities we integrate $p(x)$ over an \emph{interval}:
                $$P(X\in(a, b]) = \int _ {a}^{b} p(x) \d x$$
            \end{column}
            \begin{column}{0.45\textwidth}
                \centering
                \includegraphics[width=0.70\textwidth]{fig/prob_density_1.eps}
                \hspace*{2mm}\includegraphics[width=0.65\textwidth]{fig/prob_density_2.eps}
            \end{column}
        \end{columns}
    \end{frame}

    \begin{frame}{Cumulative Distribution Functions (CDFs)}
        \begin{boxed}
            The CDF $F(x)$ of a random variable $X$ with density $p(x)$ is given by
            $$F(x) = \int _ {-\infty}^{x} p(x) \d x = P(X \le x)$$
            Note that $F(x)$ represents an actual probability
        \end{boxed}

        Properties:
        \begin{itemize}
            \item $F(x) \rightarrow 0$ as $x\rightarrow -\infty$
            \item $F(x) \rightarrow 1$ as $x\rightarrow \infty$
            \item $F(x)$ is non-decreasing, i.e. $y > x \implies F(y) \ge F(x)$
            \item $F(x)$ is continuous (has no jumps) if $X$ is a continuous random variable
        \end{itemize}
    \end{frame}

    \begin{frame}{Characterizing Distributions}
        \framesubtitle{Expected value}
        \begin{boxed}
            \textbf{Expected value of a RV (mean)}:
            \begin{itemize}
                \item $\displaystyle \mu_X = \E X = \sum_i x_i\,p_i \quad \text{(for discrete RV $X$)}$
                \item $\displaystyle \mu_X = \E X = \int_{-\infty}^\infty x\,p(x) \d x \quad \text{(for continuous RV $X$)}$
            \end{itemize}
        \end{boxed}

        Properties:
        \begin{itemize}
            \item For any RV $X$ and numbers $a$ and $b$, we have that $\mu_{aX+b} = a\mu_X + b$
            \item For any two random variables $X$ and $Y$, we have $\mu_{X+Y} = \mu_X+\mu_Y$
        \end{itemize}

        Means characterize average behavior of a RV, which is not sufficient however:
        A profit of $CHF 10 \pm CHF 1$ per sale is very different from making $CHF 10 \pm CHF 15$ per sale.
    \end{frame}

    \conceptcheck{
        One possibility to quantify risk $r$ as expected loss. That is, by weighting the loss $l$ (e.g., in CHF)
        which some event causes with the probability $p_e$ that this event actually occurs:
        $$r := p_e\:l$$

        Now assume you want to park your car in the city of Zurich. The parking garage will charge you CHF 24.
        If you park in some quiet corner, you may get fined and pay CHF 40. The chance that you'll get caught is $2/3$.
        What's the strictly economic choice here, parking garage or taking your chances?
    }{
        The parking garage as it has the lower expected loss: $24 < \frac 2 3 \times 40 \approx 26.67$
    }

    \begin{frame}{Characterizing Distributions}
        \framesubtitle{Variance and Standard Deviation}

        \begin{boxed}
            \textbf{Variance}:
            The variance measures the variation of a RV about its mean:
            $$\sigma_X^2 = \mathrm{Var}(X) = E\left[(X-\mu_X)^2\right]$$
        \end{boxed}

        Properties:
        \begin{itemize}
            \item For any RV $X$, $\mathrm{Var}(X) \ge 0$, with $\mathrm{Var}(X) = 0$ iff $X$ constant
            \item For any RV $X$ and numbers $a$ and $b$, we have that $\mathrm{Var}(aX+b) = a^2\mathrm{Var}(X)$
            \item For any two independent RV $X$ and $Y$, we have $\mathrm{Var}(X+Y) = \mathrm{Var}(X) + \mathrm{Var}(Y)$
            \item $\displaystyle E\left[(X-\mu_X)^2\right] = E[X^2] - \mu_X^2$
        \end{itemize}

        \begin{boxed}
            \textbf{Standard Deviation (SD)}:
            The SD also measures  the variation of a RV about its mean, but has the same units as the mean:
            $$\sigma_X = \sqrt{\mathrm{Var}(X)}$$
        \end{boxed}
        Analogous properties as in the case of variance apply
    \end{frame}

    \begin{frame}{Continuous Distributions}
        \framesubtitle{A Well Behaved Example}

        Task: Compute the first two moments of the uniform distribution
        $$
            p(x) = \begin{cases}
                1 & x \in [0,1],     \\
                0 & \text{otherwise}
            \end{cases}
        $$
        \vspace*{5mm}

        Find expected value:
        $$\mu_X = \int_{-\infty}^\infty xp(x) \d x = \int_0^1 x \d x = \frac{1}{2}$$

        Find variance:
        $$\sigma_X^2 = \int_{-\infty}^\infty x^2p(x) \d x - \left(\frac{1}{2}\right)^2
            = \frac{1}{3} - \frac{1}{4} = \frac{1}{12}
        $$
        where we used above identity $\E{(X-\mu_X)^2} = \E {X^2} - \mu_X^2$
    \end{frame}

    \begin{frame}{Continuous Distributions}
        \framesubtitle{An Less Well-Behaved Example}

        \begin{columns}[onlytextwidth]
            \begin{column}{0.55\textwidth}
                Task: Compute the first two moments of a Cauchy-like distribution with the PDF given by
                $$p(x) = \frac{1}{1+x^2}$$
                \vspace*{3mm}

                Trying to compute the variance yields
                $$\sigma_X^2 = \int_{-\infty}^\infty x^2p(x) \d x
                    = \int_{-\infty}^\infty \frac{x^2}{1+x^2} \d x = \infty
                $$
                Hence the variance of $p(x)$ is undefined.\\[3mm]

                A similar problem arises for the mean, which is undefined too.
            \end{column}
            \begin{column}{0.45\textwidth}
                \begin{center}
                    \includegraphics[width=0.85\textwidth]{fig/prob_cauchy.eps}
                \end{center}
            \end{column}
        \end{columns}


    \end{frame}

    \subsection{Joint, Marginal, and Conditional Densitities}

    \begin{frame}{Joint Density}
        Similar to the univariate case, a PDF can be defined for two or more variables.

        The two RVs $X$ and $Y$ can be defined to have a \emph{joint density} $p(x, y)$ or $p_{X, Y}$.
        Then properties similar to the univariate case hold:
        \begin{itemize}
            \item $\displaystyle p(x, y) \ge 0$
            \item $\displaystyle \int _ {\mathbb{R}^2} p(x, y) \d x \d y = 1$
            \item $\displaystyle P((X, Y) \in \mathcal{D}) = \int _ {\mathcal{D}} p(x, y) \d x \d y$
        \end{itemize}

        This generalizes to any number of RVs.
    \end{frame}

    \begin{frame}{Marginal Distributions}
        When dealing with multiple variables, we often ask ``how is this one variable distributed?''.
        The distribution that answers this questions this is the \emph{marginal distribution}.

        Mathematically (often referred to as \emph{integrating out} or \emph{marginalizing out} unneeded variables):
        $$p _ X(x) = \int_{-\infty}^\infty p_{X, Y}(x, y) \d y$$

        Visually:
        \begin{center}
            \includegraphics[width=0.6\textwidth]{fig/prob_marginal_distr.eps}
        \end{center}
    \end{frame}

    \begin{frame}{Conditional Probability and Statistical  Independence}
        The conditional distribution quantifies the probability of an event occurring,
        given that another event (e.g. by assertion or by evidence) has already occurred.

        \begin{boxed}
            \textbf{Conditional Probability}:
            We define the conditional probability of ``$\ev A$ given $\ev B$'', written $P(\ev A \mid \ev B)$ to be:
            $$P(\ev A \mid \ev B) = \frac{P(\ev A, \ev B)}{P(\ev B)}
                \quad \text{or equivalently} \quad
                P(\ev A, \ev B) = P(\ev A \mid \ev B) \, P(\ev B)$$
        \end{boxed}

        Based on the above we define independence:

        \begin{boxed}
            \textbf{Statistical Independence}:
            Events $\ev A$ and $\ev B$ are defined to be \emph{statistically independent} if
            $$P(\ev A, \ev B) = P(\ev A) \, P(\ev B)$$
        \end{boxed}

        Therefore independence implies both $P(\ev A \mid \ev B) = P(\ev A)$ and $P(\ev B \mid \ev A) = P(\ev B)$
    \end{frame}

    \begin{frame}{Conditional Independence}
        % \vspace*{2mm}
        Conditional probabilities are proper probabilities $\Rightarrow$ concepts of (in)dependence also apply to them.

        \begin{boxed}
            \textbf{Conditional Independence}:
            Two RVs $A$ and $B$ are \emph{conditionally independent} given a third RV $C$ iff
            $$P(A, B \mid C) = P(A \mid C) \: P(B \mid C)$$
        \end{boxed}

        Two independent RV can become dependent when conditioning on a third RV $C$.
        Occurs when RVs $A$ and $B$ correspond to causes of RV $C$.

        Example:
        \begin{itemize}
            \item A broken bone and lung cancer might be independent in the general population
            \item Now condition on being in the hospital
            \item Conditionally, the broken bone is negatively correlated with lung cancer
            \item We say the broken bone \emph{explained away} why some person is in the hospital,
                  in turn lowering the probability that the person has lung cancer
        \end{itemize}
    \end{frame}


    \begin{frame}{Bayes' Theorem}
        By definition of conditional probabilities, $P(A, B) = P(B\mid A)\:P(A)$ and
        $P(A, B) = P(A\mid B)\:P(B)$,
        which is equivalent to
        \begin{boxed}
            \textbf{Bayes' Theorem}:
            $$P(A \mid B) = \frac{P(B\mid A)\:P(A)}{P(B)}$$
        \end{boxed}

        Please note the profound implications: It allows us to reverse the order of conditioning. If
        we know how to estimate $P(B\mid A)$, $P(A)$, and $P(B)$, then we can estimate $P(A\mid B)$.

        Example:\\
        Given  the prevalence of symptoms for a given disease,
        and the overall prevalences of the disease and symptoms, respectively,
        we can determine how likely someone is to have the disease based on their symptoms.

        Special case: Without direct access to $P(B)$, we still know that
        $$P(A \mid B) \propto P(B \mid A) P(A)$$
        and use the normalization constraint $\sum_a P(A=a \mid B) = 1$.

    \end{frame}

    \begin{frame}{Marginalization}
        Marginalization (``sum rule'') is key ingredient to make Bayes theorem work.

        It simply sums or integrates out, respectively, other variables from the joint distribution:
        $$P(B) = \sum_a P(A, B)$$
        $$p_B = \int_{-\infty}^{\infty} p(a,b)\d a$$

        The result is called marginal probability.

        In practice, where joint probabilities consist of many variables, this ``normalizer'' term
        in Bayes' theorem often becomes a computational bottleneck.
    \end{frame}

    \conceptcheck{
        It's 1963 and you participate in ``Let's Make a Deal'' (TV Show). You play this game:
        \begin{itemize}
            \item There's three doors: behind two of them there's a goat, behind the third a Ferrari
            \item First you pick a door
            \item Then the host opens one of the two other doors with a goat behind
            \item Then you can optionally switch to the other closed door or stay with your initial one
            \item Then your door will be opened and if it's the Ferrari, the car is yours
        \end{itemize}
        If you want the car, what's best to do? a) stay, b) switch, c) makes no difference\\
        Hint: Use Python or Bayes' theorem
    }{
        With loss of generality, you pick door 1 and the host opens door 2 or 3. Let $D := \text{``number of the door with the car behind''} \in \{1,2,3\}$,
        and likewise $H := \text{number of door which host opens}$. At the end of the game there's two options, $d=1$ and $d=3$:\\
        $$P(D = d \mid H = 2) = \frac{P(H = 2 \mid D = d)\, P(D = d)}{P(H = 2)}$$
        $$
            P(D = 1 \mid H = 2) = \frac{1/2 \times 1/3}{1/3 + 1/6} = \frac 1 3
            \quad \text \quad
            P(D = 3 \mid H = 2) = \frac{1 \times 1/3}{1/3 + 1/6} = \frac 2 3
        $$
    }

    \begin{frame}{Covariance}
        \begin{boxed}
            \textbf{Covariance}:
            Covariance measures the degree that two random variables fluctuate together:
            $$\sigma_{XY} = \mathrm{Cov}(X, Y) = E[XY] - E[X]E[Y]$$
        \end{boxed}

        Spelled out this means for the discrete and continuous case, respectively,
        $$\sigma_{XY} = \sum_{i, j} (x_i - \mu_X) (y_j-\mu_Y) p_{ij} \quad \text{and} \quad
            \sigma_{XY} = \int _ {\mathbb{R}^2} (x - \mu_X) (y - \mu_y) p(x, y) \d x \d y$$

        Visually:
        \vspace*{-5mm}
        \begin{center}
            \includegraphics[width=0.6\textwidth]{fig/prob_cov.eps}
        \end{center}
    \end{frame}

    \begin{frame}
        Summarizing the main properties of covariance:
        \begin{itemize}
            \item For any RV $X$, $\mathrm{Cov}(X, X) = \mathrm{Var}(X)$
            \item For any RVs $X$ and $Y$ and numbers $a$ and $b$,
                  $\mathrm{Cov}(aX+b, Y) = \mathrm{Cov}(X, aY+b) = a\mathrm{Cov}(X, Y)$
            \item If $X$ and $Y$ are independent then $\mathrm{Cov}(X, Y) = 0$
        \end{itemize}

    \end{frame}

    \begin{frame}{Correlation}
        \vspace*{-2mm}
        \begin{boxed}
            \textbf{Correlation coefficient}:
            Define the \emph{correlation coefficient} to be
            $$\rho(X, Y) = \frac{\mathrm{Cov}(X, Y)}{\sigma_{X}\sigma_{Y}}$$
        \end{boxed}

        \vspace*{-1mm}
        Properties:
        \vspace*{-1mm}
        \begin{itemize}
            \item For any RV $X$, $\rho(X, X) = 1$
            \item For any RVs $X$ and $Y$ and numbers $a$ and $b$, $\rho(aX+b, Y) = \rho(X, aY+b) = \rho(X, Y)$
            \item If $X$ and $Y$ are independent with non-zero variance then $\rho(X, Y) = 0$
        \end{itemize}

        Visually:
        \vspace*{-5mm}
        \begin{center}
            \includegraphics[width=0.6\textwidth]{fig/prob_corr.eps}
        \end{center}
    \end{frame}

    \begin{frame}{}
        \vspace*{2mm}
        \textbf{Example}\\[3mm]

        Let $X$ be any random variable, and $Y=aX+b$ as linear deterministic function of $X$.

        Then, one can compute that
        $$\sigma_{Y} = \sigma_{aX+b} = |a|\,\sigma_{X}$$

        $$\mathrm{Cov}(X, Y) = \mathrm{Cov}(X, aX+b) = a\mathrm{Cov}(X, X) = a\mathrm{Var}(X),$$

        and by definition

        $$\rho(X, Y) = \frac{a\mathrm{Var}(X)}{|a|\sigma_{X}^2} = \frac{a}{|a|} = \mathrm{sign}(a)$$

        Thus the correlation is $+1$ for any $a > 0$, and $-1$ for any $a < 0$.

        This shows that correlation
        measures the \emph{degree and directionality} of variation of the two RV, not the scale
        that the variation takes.
    \end{frame}

    % TODO
    % \begin{frame}{title}
    %     Show example: \url{http://d2l.ai/chapter_preliminaries/probability.html#an-example}
    %  \end{frame}

    \begin{frame}{Maximum Likelihood}

        Maximum Likelihood:
        \begin{itemize}
            \item A very common way of thinking in machine learning. Really it is a ``point of view''
            \item Given a probabilistic model with unknown parameters,
                  the most likely parameters are those under which the chance
                  that the model generated the data are maximized
        \end{itemize}

        Technically, given our data $X$, we want
        $$\hat{\boldsymbol{\theta}} = \arg \max_{\boldsymbol{\theta}} P(\boldsymbol{\theta}\mid X)$$
        where $\boldsymbol{\theta} = \mqty(\theta_1, \ldots, \theta_k)^\top$ the parameter vector
        of interest.

        Using Bayes' theorem the Maximum Likelihood Estimator (MLE) for $\boldsymbol{\theta}$ is
        $$\hat{\boldsymbol{\theta}}
            = \arg \max_{\boldsymbol{\theta}} \frac{P(X \mid \boldsymbol{\theta})P(\boldsymbol{\theta})}{P(X)}
            = \arg \max_{\boldsymbol{\theta}} P(X \mid \boldsymbol{\theta})P(\boldsymbol{\theta})$$
        since $P(X)$ is just scaling the objective w.r.t. $\boldsymbol{\theta}$.

        Exact same arguments apply to the densities of the continuous case.
    \end{frame}

    \begin{frame}{Maximum Likelihood: An Example}
        \begin{columns}[onlytextwidth]
            \begin{column}{0.7\textwidth}
                Setting:
                \begin{itemize}
                    \item Given a single, unknown parameter $\theta$: Probability that a coin flip is heads
                    \item Assume ``uniformative prior'' $\Rightarrow$ drop $P(\theta)$ from Bayes' theorem
                    \item Data: ``HHHTHTTHHHHHT'', i.e. $n_H = 9$ heads, $n_T = 4$ tails
                \end{itemize}
                \vspace*{3mm}

                Question: What is the probability that the coin comes up heads?\\
            \end{column}
            \begin{column}{0.3\textwidth}
                \begin{center}
                    \includegraphics[width=\textwidth]{fig/prob_mle.eps}
                \end{center}
            \end{column}
        \end{columns}

        Answer: Assume samples independent $\Rightarrow$ probabilities multiply. Likelihood becomes
        $$P(X \mid \theta) = \theta^{n_H}(1-\theta)^{n_T} = P(X \mid \theta) = \theta^9(1-\theta)^4$$

        Then maximize the likelihood: The first-order condition
        $$
            \frac{d}{d\theta} P(X \mid \theta)
            = \frac{d}{d\theta} \theta^9(1-\theta)^4
            = 9\theta^8(1-\theta)^4 - 4\theta^9(1-\theta)^3
            = \theta^8(1-\theta)^3(9-13\theta)
            = 0
        $$
        has three solutions: $0$, $1$ and $9/13$. First two are minima since
        $P(X\mid \theta=0) = P(X\mid \theta=1) = 0$. Last solution assigns non-zero probability
        to our data, hence the maximum likelihood estimate is $\hat \theta = 9/13.$\\
        {\tiny Note that despite everyone would correctly guess $9/13$, MLE provides
        a structured approach which equally applies to vastly more complex situations}

    \end{frame}

    \begin{frame}{Numerical Optimization and the Negative Log-Likelihood}
        In practice, we may have billions of parameters and examples.
        \begin{itemize}
            \item Independence assumption results in long product of small values, say, e.g.
                  $\prod_{i=1}^{10^9} \frac12$, something far below machine precision
            \item However, the $\log$ of this product is a sum,
                  $-\sum_{i=0}^{10^9} \log_{10} 2 = -10^9 \times \log_{10} 2 \approx -301029995.664$,
                  which fits into a 32-bit floating point value
        \end{itemize}

        Hence should work with log-likelihood,
        $$\log P(X \mid \boldsymbol{\theta})$$
        which will not affect the result since $\log$ is monotonically increasing.

        Often working with loss functions, to be \emph{minimized}, hence use negative log-likelihood:
        $$-\log P(X \mid \boldsymbol{\theta})$$

        Example: For the the coin flipping problem from before, this means minimizing
        $$-\log(P(X \mid \boldsymbol{\theta})) = -\log(\theta^{n_H}(1-\theta)^{n_T}) = -n_H\log(\theta) - n_T\log(1-\theta)$$
        which again leads to $\hat{\theta} = n_H/(n_H + n_T)$
    \end{frame}

    \subsection{Basic Distributions}

    \begin{frame}{Distributions}
        \framesubtitle{Bernoulli}
        Simple RV encoding e.g. a coin flip which comes up $1$ (heads)
        with probability $p$ and $0$ (tails) with probability $1-p$.

        \vspace*{3mm}
        \begin{columns}[onlytextwidth]
            \begin{column}{0.5\textwidth}
                \begin{boxed}
                    \textbf{Bernoulli Distribution}:
                    $$X \sim \mathrm{Bernoulli}(p) $$

                    Cumulative distribution function:
                    $$F(x) = \begin{cases} 0 & x < 0, \\ 1-p & 0 \le x < 1, \\ 1 & x \geq 1 . \end{cases}$$

                    Moments:
                    \begin{itemize}
                        \item  $\mu_X = p$
                        \item $\sigma_X^2 = p(1-p)$
                    \end{itemize}
                \end{boxed}
            \end{column}
            \begin{column}{0.5\textwidth}
                \begin{center}
                    \includegraphics[width=0.8\textwidth]{fig/prob_dist_bernoulli.eps}
                \end{center}
            \end{column}
        \end{columns}
    \end{frame}

    \begin{frame}{Distributions}
        \framesubtitle{Discrete Uniform}
        Every value out of given set is equally likely has a probability of $p_i = \frac{1}{n}$ to occur

        \vspace*{3mm}
        \begin{columns}[onlytextwidth]
            \begin{column}{0.5\textwidth}
                \begin{boxed}
                    \textbf{Discrete Uniform Distribution}:
                    Assume the set of possible outcomes is $\{1, 2, \ldots, n\}$.
                    $$X \sim U(n)$$

                    Cumulative distribution function:
                    $$F(x) = \begin{cases} 0 & x < 1, \\ \frac{k}{n} & k \le x < k+1 \text{ with } 1 \le k < n, \\ 1 & x >= n . \end{cases}$$

                    Moments:
                    \begin{itemize}
                        \item $\mu_X = \frac{1+n}{2}$
                        \item $\sigma_X^2 = \frac{n^2-1}{12}$
                    \end{itemize}
                \end{boxed}
            \end{column}
            \begin{column}{0.5\textwidth}
                \begin{center}
                    \includegraphics[width=0.8\textwidth]{fig/prob_dist_uniform.eps}
                \end{center}
            \end{column}
        \end{columns}
    \end{frame}

    \begin{frame}{Distributions}
        \framesubtitle{Continuous Uniform}
        \vspace*{-1mm}
        Similar to discrete uniform distribution: Starting from $U(n)$ with $n \rightarrow \infty$
        and linearly scaling it to the interval $[a,b]$ will approach a continuous random variable that just
        picks an arbitrary value in $[a,b]$ with equal probability

        \begin{columns}[onlytextwidth]
            \begin{column}{0.5\textwidth}
                \begin{boxed}
                    \textbf{Continuous Uniform Distribution}:
                    $$X \sim U(a, b)$$

                    $$\text{PDF:}\qquad p(x) = \begin{cases} \frac{1}{b-a} & x \in [a, b] \\ 0 & x \not\in [a, b]\end{cases}$$

                    $$\text{CDF:}\qquad F(x) = \begin{cases} 0 & x < a \\ \frac{x-a}{b-a} & x \in [a, b] \\ 1 & x >= b \end{cases}$$

                    $$\text{Moments:}\qquad \displaystyle \mu_X = \frac{a+b}{2} \quad \sigma_X^2 = \frac{(b-a)^2}{12}$$
                \end{boxed}
            \end{column}
            \begin{column}{0.5\textwidth}
                \begin{center}
                    \includegraphics[width=0.55\textwidth]{fig/prob_cont_unif_pdf.eps}
                    \includegraphics[width=0.55\textwidth]{fig/prob_cont_unif_cdf.eps}
                \end{center}
            \end{column}
        \end{columns}
    \end{frame}

    \begin{frame}{Distributions}
        \framesubtitle{Binomial}
        Given $n$ independent experiments, each with success probability $p$. Interested in
        number of successes $X$:
        $$X = \sum_{i=1}^n X_i \quad \text{with} \quad X_i \sim \mathrm{Bernoulli}(p)
            \qquad \Longleftrightarrow \qquad X \sim \mathrm{Binomial}(n, p)
        $$

        \vspace*{-2mm}
        \begin{boxed}
            \textbf{Binomial Distribution}:
            $$
                \text{CDF:}\qquad F(x) = \begin{cases}
                    0                                          & x < 0                                   \\
                    \sum_{m \le k} \binom{n}{m} p^m(1-p)^{n-m} & k \le x < k+1 \text{ with } 0 \le k < n \\
                    1                                          & x >= n
                \end{cases}
            $$

            $$\text{Moments:}\qquad \displaystyle \mu_X = np \quad \sigma_X^2 = np(1-p)$$
        \end{boxed}
        \vspace*{-3mm}
        \begin{center}
            \includegraphics[width=0.3\textwidth]{fig/prob_binomial_pdf.eps}
            \qquad
            \includegraphics[width=0.3\textwidth]{fig/prob_binomial_cdf.eps}
        \end{center}

    \end{frame}

    \begin{frame}{Distributions}
        \framesubtitle{{Poisson}}
        \vspace*{-8mm}
        \begin{columns}[onlytextwidth]
            \begin{column}{0.57\textwidth}
                Model for arrival rates, e.g. number of buses at a given stop per
                time interval. Call this rate $\lambda$ (with units [1/s])
                \vspace*{2mm}

                \begin{boxed}
                    \textbf{Poisson Distribution}:
                    $$X \sim U(a, b)$$

                    $$\text{PMF:}\qquad p_k = \frac{\lambda^ke^{-\lambda}}{k!}$$

                    $$\text{CDF:}\qquad F(x) = \begin{cases} 0 & x < 0, \\ e^{-\lambda}\sum_{m = 0}^k \frac{\lambda^m}{m!} & k \le x < k+1 \text{ with } 0 \le k \end{cases}$$

                    $$\text{Moments:}\qquad \displaystyle \mu_X = \lambda \quad \sigma_X^2 = \lambda$$
                \end{boxed}
            \end{column}
            \begin{column}{0.43\textwidth}
                \begin{center}
                    \includegraphics[width=0.8\textwidth]{fig/prob_poisson_pmf.eps}
                    \includegraphics[width=0.8\textwidth]{fig/prob_poisson_cdf.eps}
                \end{center}
            \end{column}
        \end{columns}
    \end{frame}

    \begin{frame}{Distributions}
        \framesubtitle{Gaussian}

        \begin{columns}[onlytextwidth]
            \begin{column}{0.55\textwidth}
                The central distribution in probability theory with many unique properties because of
                the \emph{central limit theorem}.\\[1mm]

                Informally, given $N$ arbitrarily distributed, i.i.d. RVs $X_i$ and
                $$
                    X^{(N)} = \sum_{i=1}^N X_i
                $$
                their sum. Then, under mild conditions,
                $$
                    \frac{X^{(N)} - \mu_{X^{(N)}}}{\sigma_{X^{(N)}}} \quad \sim \quad \text{approximately Gaussian}
                $$

                \begin{boxed}
                    \textbf{Gaussian Distribution}:
                    $$X \sim \mathcal{N}(\mu, \sigma^2)$$
                    \vspace*{-2mm}
                    $$\text{PDF:}\qquad p_X(x) = \frac{1}{\sqrt{2\pi}\sigma}e^{-\frac{(x-\mu)^2}{2\sigma^2}}$$
                    $$\text{Moments:}\qquad \displaystyle \mu_X = \mu \quad \sigma_X^2 = \sigma^2$$
                \end{boxed}

            \end{column}
            \begin{column}{0.45\textwidth}
                \begin{center}
                    \includegraphics[width=0.73\textwidth]{fig/prob_norm_pdf.eps}
                    \qquad
                    \includegraphics[width=0.73\textwidth]{fig/prob_norm_cdf.eps}
                \end{center}

            \end{column}
        \end{columns}
    \end{frame}


    \begin{frame}{Distributions: Summary}
        Things to remember about distributions:
        \begin{itemize}
            \item Bernoulli RVs can be used to model events with a yes/no outcome
            \item Discrete uniform distributions model selects from a finite set of possibilities
            \item Continuous uniform distributions select from an interval
            \item Binomial distributions model a series of Bernoulli RVs, and count the number of successes
            \item Poisson RVs model the arrival of rare events
            \item Gaussian RVs model the result of adding a large number of independent RVs together
            \item All the above distributions belong to the exponential family
        \end{itemize}
    \end{frame}

    % \begin{frame}{Statistics}
    %     probably should discuss this as well,
    %     \url{http://d2l.ai/chapter_appendix-mathematics-for-deep-learning/statistics.html}
    % \end{frame}
}

% \include{pt4}
% \include{pt5}

\begin{frame}{License}
    \small
    This work is based on the following works:
    \begin{itemize}
        \item Jim Hefferson, \emph{Linear Algebra}, 4th Edition, freely available at \url{https://hefferon.net/linearalgebra/}
        \item Wikipedia, The Free Encyclopedia, \url{https://www.wikipedia.org/}
        \item D2L.ai: Interactive Deep Learning Book with Multi-Framework Code, Math, and Discussions, freely available at \url{https://github.com/d2l-ai/d2l-en}
    \end{itemize}

    The above three works are based on Creative Commons Attribution-ShareAlike (CC BY-SA)-compatible licenses. Hence this work is again released under
    the same permissive Creative Commons Zero v1.0 Universal License. Sources are available at \url{https://github.com/openteaching/ml-math-refresher}
\end{frame}

\end{document}
